% ============================================================
\chapter{Levi-Civita Connection}
\label{ch:levi_civita}
% ============================================================

On a curved surface, the Euclidean notion of ``straight line'' loses its meaning. But the idea of ``the straightest possible line'' survives: it is the geodesic. To define it, one needs to know how to differentiate vector fields along curves on the manifold---and this is where the Levi-Civita connection enters, named after the Italian mathematician Tullio Levi-Civita who, together with Gregorio Ricci-Curbastro, developed tensor calculus at the turn of the twentieth century. The fundamental result of this chapter is as surprising as it is elegant: there exists a \emph{unique} way to differentiate that is compatible with the metric and torsion-free. This uniqueness is the fundamental theorem of Riemannian geometry.

\section{The fundamental theorem of Riemannian geometry}

\begin{theorem}[Fundamental theorem]
\label{thm:levi_civita}
Let $(M, g)$ be a Riemannian manifold. There exists a unique connection $\nabla$
on $TM$ satisfying:
\begin{enumerate}
    \item \textbf{Metric compatibility:} $\nabla g = 0$, i.e.\
    $X\inner{Y}{Z} = \inner{\nabla_X Y}{Z} + \inner{Y}{\nabla_X Z}$.
    \item \textbf{Torsion-freeness:} $\nabla_X Y - \nabla_Y X = [X, Y]$.
\end{enumerate}
This connection is called the \emph{Levi-Civita connection} of $(M, g)$.
\end{theorem}

\begin{proof}
\textbf{Uniqueness.} Assume $\nabla$ satisfies (1) and (2). Writing compatibility
for three cyclic permutations:
\begin{align*}
    X\inner{Y}{Z} &= \inner{\nabla_X Y}{Z} + \inner{Y}{\nabla_X Z}, \\
    Y\inner{Z}{X} &= \inner{\nabla_Y Z}{X} + \inner{Z}{\nabla_Y X}, \\
    Z\inner{X}{Y} &= \inner{\nabla_Z X}{Y} + \inner{X}{\nabla_Z Y}.
\end{align*}
Taking the combination $(1) + (2) - (3)$ and using torsion-freeness
($\nabla_X Y - \nabla_Y X = [X,Y]$), we obtain the \emph{Koszul formula}:
\begin{align}
    2\inner{\nabla_X Y}{Z} &= X\inner{Y}{Z} + Y\inner{Z}{X} - Z\inner{X}{Y}
    \notag \\
    &\quad + \inner{[X,Y]}{Z} - \inner{[Y,Z]}{X} + \inner{[Z,X]}{Y}.
    \label{eq:koszul}
\end{align}
The right-hand side is entirely determined by $g$ and the Lie brackets, so
$\nabla_X Y$ is uniquely determined by the nondegeneracy of $g$.

\textbf{Existence.} Define $\nabla_X Y$ by the Koszul formula~\eqref{eq:koszul}.
One verifies that this defines a connection (linearity in $X$, Leibniz rule
in $Y$), that it is metric-compatible, and torsion-free. The verifications are
direct but computational.
\end{proof}

\begin{keyformulas}[title=\textbf{Koszul formula}]
\[
    2\inner{\nabla_X Y}{Z} = X\inner{Y}{Z} + Y\inner{X}{Z} - Z\inner{X}{Y}
    + \inner{[X,Y]}{Z} - \inner{[X,Z]}{Y} - \inner{[Y,Z]}{X}.
\]
This is \emph{the} fundamental formula of Riemannian geometry.
\end{keyformulas}

\section{Christoffel symbols of the Levi-Civita connection}

\begin{proposition}[Christoffel symbol formula]
In local coordinates $(x^1, \ldots, x^n)$, the Christoffel symbols of the
Levi-Civita connection are:
\[
    \Gamma^k_{ij} = \frac{1}{2} g^{kl}
    \left(\frac{\partial g_{jl}}{\partial x^i}
    + \frac{\partial g_{il}}{\partial x^j}
    - \frac{\partial g_{ij}}{\partial x^l}\right).
\]
\end{proposition}

\begin{proof}
Applying the Koszul formula to coordinate vector fields $\partial_i$, $\partial_j$,
$\partial_l$ (whose Lie brackets vanish):
\[
    2\inner{\nabla_{\partial_i}\partial_j}{\partial_l}
    = \partial_i g_{jl} + \partial_j g_{il} - \partial_l g_{ij}.
\]
Hence $2\, g_{kl}\, \Gamma^k_{ij} = \partial_i g_{jl} + \partial_j g_{il} - \partial_l g_{ij}$,
and multiplying by $g^{lm}/2$ yields the formula.
\end{proof}

\begin{remark}
The symmetry $\Gamma^k_{ij} = \Gamma^k_{ji}$ is immediate, reflecting
torsion-freeness: $[\partial_i, \partial_j] = 0$ implies
$\nabla_{\partial_i}\partial_j = \nabla_{\partial_j}\partial_i$.
\end{remark}

\section{Computations on fundamental examples}

\begin{example}[Sphere $S^2$ in spherical coordinates]
With $g = d\theta^2 + \sin^2\!\theta\, d\phi^2$:
\[
    g_{11} = 1, \quad g_{22} = \sin^2\!\theta, \quad g_{12} = 0.
\]
The nonzero Christoffel symbols are:
\begin{align*}
    \Gamma^1_{22} &= -\sin\theta\cos\theta, \\
    \Gamma^2_{12} = \Gamma^2_{21} &= \cot\theta.
\end{align*}
All others vanish. Verification:
$\Gamma^1_{22} = \frac{1}{2}g^{11}(-\partial_1 g_{22}) = -\frac{1}{2}\cdot 2\sin\theta\cos\theta
= -\sin\theta\cos\theta$.
\end{example}

\begin{example}[Hyperbolic space $\mathbb{H}^2$ (half-plane)]
With $g = \frac{dx^2 + dy^2}{y^2}$:
\[
    g_{11} = g_{22} = \frac{1}{y^2}, \quad g_{12} = 0,
    \quad g^{11} = g^{22} = y^2, \quad g^{12} = 0.
\]
The nonzero symbols are:
\begin{align*}
    \Gamma^1_{12} = \Gamma^1_{21} &= -\frac{1}{y}, \\
    \Gamma^2_{11} &= \frac{1}{y}, \\
    \Gamma^2_{22} &= -\frac{1}{y}.
\end{align*}
\end{example}

\begin{example}[Surface of revolution]
Let $S$ be the surface of revolution parametrized by
$(u, v) \mapsto (f(u)\cos v, f(u)\sin v, h(u))$ with $f > 0$.
The induced metric is $g = (f'^2 + h'^2)\, du^2 + f^2\, dv^2$.

If the generating curve is arc-length parametrized ($f'^2 + h'^2 = 1$),
then $g = du^2 + f(u)^2\, dv^2$ and the nonzero symbols are:
\[
    \Gamma^1_{22} = -f f', \qquad
    \Gamma^2_{12} = \Gamma^2_{21} = \frac{f'}{f}.
\]
\end{example}

\begin{example}[Lie group with bi-invariant metric]
\label{ex:bi_inv_christoffel}
On a Lie group $G$ with bi-invariant metric, for left-invariant fields
$X, Y \in \mathfrak{g}$, the Koszul formula gives:
\[
    \nabla_X Y = \frac{1}{2}[X, Y].
\]
Indeed, bi-invariance gives $\inner{[X,Y]}{Z} + \inner{Y}{[X,Z]} = 0$
(ad-antisymmetry), and the terms $X\inner{Y}{Z}$ etc.\ vanish since the
fields and the metric are left-invariant. The Koszul formula reduces to
$2\inner{\nabla_X Y}{Z} = \inner{[X,Y]}{Z} + \inner{[Z,X]}{Y} - \inner{[Y,Z]}{X}
= \inner{[X,Y]}{Z}$.
\end{example}

\section{Gradient, divergence, and Laplacian}

\begin{definition}[Gradient]
The \emph{gradient} of $f \in C^\infty(M)$ is the unique vector field
$\operatorname{grad} f = (\nabla f)^\sharp$ such that:
\[
    g(\operatorname{grad} f, X) = Xf = df(X), \quad \forall\, X \in \mathfrak{X}(M).
\]
In coordinates: $(\operatorname{grad} f)^i = g^{ij} \frac{\partial f}{\partial x^j}$.
\end{definition}

\begin{definition}[Divergence]
The \emph{divergence} of a vector field $X \in \mathfrak{X}(M)$ is:
\[
    \operatorname{div} X = \tr(\nabla X) = \nabla_i X^i
    = \frac{1}{\sqrt{\det g}} \frac{\partial}{\partial x^i}
    \left(\sqrt{\det g}\, X^i\right).
\]
\end{definition}

\begin{definition}[Laplace--Beltrami operator]
The \emph{Laplacian} of a function $f$ is:
\[
    \Delta f = \operatorname{div}(\operatorname{grad} f)
    = \frac{1}{\sqrt{\det g}} \frac{\partial}{\partial x^i}
    \left(\sqrt{\det g}\, g^{ij} \frac{\partial f}{\partial x^j}\right).
\]
\end{definition}

\begin{example}[Laplacian on $S^2$]
In spherical coordinates, $\sqrt{\det g} = \sin\theta$, so:
\[
    \Delta_{S^2} f = \frac{1}{\sin\theta}\frac{\partial}{\partial\theta}
    \left(\sin\theta\, \frac{\partial f}{\partial\theta}\right)
    + \frac{1}{\sin^2\!\theta}\frac{\partial^2 f}{\partial\phi^2}.
\]
\end{example}

\begin{example}[Laplacian on $\mathbb{H}^2$]
With $g = y^{-2}(dx^2 + dy^2)$:
\[
    \Delta_{\mathbb{H}^2} f = y^2\left(\frac{\partial^2 f}{\partial x^2}
    + \frac{\partial^2 f}{\partial y^2}\right).
\]
\end{example}

\section{Covariant differentiation along curves}

\begin{definition}
Let $\gamma : [a,b] \to M$ be a smooth curve and $V$ a vector field along
$\gamma$. The \emph{covariant derivative} of $V$ along $\gamma$ is:
\[
    \frac{DV}{dt} = \nabla_{\dot\gamma} V.
\]
In coordinates: $\displaystyle\left(\frac{DV}{dt}\right)^k
= \frac{dV^k}{dt} + \Gamma^k_{ij}(\gamma(t))\, \dot\gamma^i(t)\, V^j(t)$.
\end{definition}

\begin{proposition}[Leibniz rule along curves]
If $V$ and $W$ are vector fields along $\gamma$ and $\nabla$ is metric-compatible:
\[
    \frac{d}{dt}\inner{V}{W} = \inner{\frac{DV}{dt}}{W} + \inner{V}{\frac{DW}{dt}}.
\]
\end{proposition}

\section{Covariant acceleration and geodesics}

\begin{definition}[Covariant acceleration]
The \emph{covariant acceleration} of a curve $\gamma$ is:
\[
    \frac{D\dot\gamma}{dt} = \nabla_{\dot\gamma}\dot\gamma.
\]
In coordinates: $\displaystyle\left(\frac{D\dot\gamma}{dt}\right)^k
= \ddot\gamma^k + \Gamma^k_{ij}\, \dot\gamma^i\, \dot\gamma^j$.
\end{definition}

\begin{definition}[Geodesic]
A curve $\gamma$ is a \emph{geodesic} if its covariant acceleration vanishes:
\[
    \frac{D\dot\gamma}{dt} = \nabla_{\dot\gamma}\dot\gamma = 0.
\]
\end{definition}

The geodesic equation in coordinates is the second-order ODE system:
\[
    \ddot\gamma^k + \Gamma^k_{ij}(\gamma)\, \dot\gamma^i\, \dot\gamma^j = 0,
    \quad k = 1, \ldots, n.
\]

\begin{keyformulas}[title=\textbf{Geodesic equation}]
\[
    \ddot\gamma^k + \Gamma^k_{ij}\, \dot\gamma^i\, \dot\gamma^j = 0.
\]
A system of $n$ second-order ODEs, equivalent to $2n$ first-order ODEs on $TM$.
\end{keyformulas}

\section{Levi-Civita connection of products and warped products}

\begin{proposition}[Riemannian product]
On $(M_1 \times M_2, g_1 \oplus g_2)$, if $X_1, Y_1 \in \mathfrak{X}(M_1)$ and
$X_2, Y_2 \in \mathfrak{X}(M_2)$, then:
\[
    \nabla_{(X_1, X_2)}(Y_1, Y_2) = (\nabla^1_{X_1} Y_1, \nabla^2_{X_2} Y_2),
\]
where $\nabla^1$ and $\nabla^2$ are the Levi-Civita connections of $(M_1, g_1)$
and $(M_2, g_2)$ respectively.
\end{proposition}

\begin{proposition}[Warped product]
\label{prop:warped_LC}
On the warped product $B \times_f F$ with $g = g_B + f^2 g_F$, the Levi-Civita
connection satisfies, for $X, Y \in \mathfrak{X}(B)$ and $V, W \in \mathfrak{X}(F)$:
\begin{align*}
    \nabla_X Y &= \nabla^B_X Y, \\
    \nabla_X V = \nabla_V X &= \frac{Xf}{f}\, V, \\
    \nabla_V W &= \nabla^F_V W - \frac{\inner{V}{W}_F}{f}\, \operatorname{grad}_B f.
\end{align*}
\end{proposition}

\section{Exercises}

\begin{exercise}
Compute all Christoffel symbols for the metric
$g = dr^2 + r^2(d\theta^2 + \sin^2\!\theta\, d\phi^2)$ of $\R^3$ in spherical
coordinates.
\end{exercise}

\begin{exercise}
Verify that great circles on $S^2$ are geodesics using the Christoffel symbols
computed above.
\end{exercise}

\begin{exercise}
Show that on a Lie group with bi-invariant metric, the geodesics through the
identity $e$ are exactly the one-parameter subgroups $t \mapsto \exp(tX)$.
\end{exercise}

\begin{exercise}
Compute the Laplacian of $f(x,y) = \ln y$ on $\mathbb{H}^2$ with metric
$g = \frac{dx^2 + dy^2}{y^2}$.
\end{exercise}

\begin{exercise}
Let $(M, g)$ be Riemannian and $\tilde{g} = e^{2\varphi} g$ a conformal change.
Show that the Christoffel symbols are related by:
\[
    \tilde{\Gamma}^k_{ij} = \Gamma^k_{ij} + \delta^k_i \partial_j \varphi
    + \delta^k_j \partial_i \varphi - g_{ij} g^{kl} \partial_l \varphi.
\]
\end{exercise}

\begin{exercise}
On the warped product $\R \times_{\cosh t} \R^{n-1}$ with
$g = dt^2 + \cosh^2(t)\, g_{\R^{n-1}}$, compute the Christoffel symbols
and verify that the curves $t \mapsto (t, x_0)$ are geodesics.
\end{exercise}
